\section{第一阶段证据与验收：用可观察现象校准理解}

第一阶段的目标是让硬件名词落到可观察现象，而不是让读者背出更长清单。真正理解 CPU、MMU、DMA、中断、cache 和设备队列，应该能在系统上找到证据：日志、计数器、地址映射、页错误、块 I/O、网络队列、调度事件和错误报告。

\subsection{观察启动链：哪些硬件事实被交给内核}

启动日志是硬件到 OS 的第一份证据。它会显示内存地图、CPU 特性、NUMA 拓扑、时钟源、中断控制器、PCIe 设备、IOMMU、存储设备和根文件系统。读启动日志时，不要只找错误，要把它当成硬件契约清单。

\begin{lstlisting}
evidence to collect on a Linux system:
  dmesg
  cat /proc/cpuinfo
  cat /proc/iomem
  cat /proc/interrupts
  lspci -vv
  lsblk
  numactl --hardware
\end{lstlisting}

\begin{longtable}{p{0.22\linewidth}p{0.34\linewidth}p{0.32\linewidth}}
\toprule
证据 & 说明什么 & 关联章节 \\
\midrule
\code{/proc/iomem} & RAM、保留区、MMIO 窗口 & 物理内存地图、设备地址 \\
\code{/proc/interrupts} & IRQ 分配和计数 & 中断控制器、设备完成 \\
\code{lspci -vv} & BAR、MSI-X、AER、IOMMU group & PCIe、驱动资源 \\
\code{dmesg} & probe、错误、时钟、IOMMU、NUMA & 启动链、设备模型 \\
\code{numactl --hardware} & CPU 与内存 node 距离 & NUMA、页放置 \\
\bottomrule
\end{longtable}

这些命令是原理卷的校准工具，不替代实践卷实验。书里说“MMIO 是地址空间的一部分”，\code{/proc/iomem} 应该能看到设备窗口；书里说“中断有亲和性和计数”，\code{/proc/interrupts} 应该能看到每个 CPU 上的计数分布。

\subsection{观察系统调用和特权边界}

\code{strace} 能看到用户态请求了哪些系统调用，但看不到内核内部所有路径。它适合验证“库函数不是系统调用本身”和“系统调用有返回值与错误码”。例如一个 \code{read} 可能返回短读，一个 \code{open} 可能因权限失败，一个 \code{mmap} 可能成功但物理页尚未分配。

\begin{lstlisting}
questions while reading strace:
  Which calls cross the user/kernel boundary?
  Which file descriptors are created, duplicated, or closed?
  Which calls return EAGAIN, EINTR, EFAULT, or EPERM?
  Does mmap only create virtual ranges, or also force page allocation?
\end{lstlisting}

特权边界的负面证据同样重要。用户态传入坏指针时，内核应返回 \code{EFAULT} 或终止进程，而不是崩溃；无权限打开文件应失败且不产生副作用；被 seccomp 拒绝的系统调用不应进入危险路径。

\subsection{观察虚拟内存：VMA、RSS 和 page fault}

进程的虚拟地址空间可以从 \code{/proc/<pid>/maps} 和 \code{/proc/<pid>/smaps} 观察。maps 显示 VMA，smaps 显示更细的 RSS、dirty、shared/private 等统计。page fault 可以用 \code{perf stat} 或进程资源统计观察。

\begin{lstlisting}
memory evidence:
  cat /proc/<pid>/maps
  cat /proc/<pid>/smaps
  perf stat -e page-faults,minor-faults,major-faults <command>
  pmap -x <pid>
\end{lstlisting}

\begin{longtable}{p{0.22\linewidth}p{0.34\linewidth}p{0.32\linewidth}}
\toprule
现象 & 可能原因 & 需要区分 \\
\midrule
VIRT 很大 RSS 很小 & VMA 已建，页未实际分配 & 懒分配、文件映射、保留地址空间 \\
minor fault 多 & 分配零页或建立 PTE & 不一定访问磁盘 \\
major fault 多 & 需要存储 I/O & 文件页或 swap 读入 \\
shared 增加 & 共享库或共享映射 & 不等于进程独占内存增加 \\
dirty 增加 & 写过但未回写 & page cache 或匿名页 \\
\bottomrule
\end{longtable}

这组证据对应的核心理解是：虚拟地址、VMA、PTE、物理页、RSS 不是同一层概念。一个地址可见，不代表物理页已分配；一个页在 RSS 中，不代表它独占；一个 page fault 发生，不代表程序错误。

\subsection{观察调度：on-CPU 与 off-CPU 分开}

调度问题最容易被误判。请求慢可能是 CPU 算得慢，也可能是 runnable 但排队，也可能是 sleeping 等 I/O 或锁。证据上要分 on-CPU、runnable wait 和 blocked wait。

\begin{lstlisting}
scheduler evidence:
  perf sched record/report
  trace-cmd record -e sched:*
  cat /proc/<pid>/sched
  cat /proc/schedstat
  pidstat -w
\end{lstlisting}

若上下文切换次数非常高，可能是线程过多、锁竞争、I/O 小请求过多或时间片太短。若 runnable wait 高，说明任务想运行但拿不到 CPU。若 CPU 使用率低但延迟高，可能大量时间在 I/O、锁或内存回收中睡眠。

\subsection{观察设备和 I/O：队列、completion、错误}

块设备和网络设备都有队列。吞吐不只由设备峰值决定，还由队列深度、请求大小、合并、完成延迟、中断亲和性和 CPU 处理能力共同决定。

\begin{lstlisting}
I/O evidence:
  iostat -x
  lsblk -o NAME,MODEL,ROTA,SIZE
  cat /sys/block/<dev>/queue/*
  ethtool -S <iface>
  cat /proc/interrupts
  dmesg | grep -i -E "nvme|pcie|iommu|dma|reset|timeout"
\end{lstlisting}

\begin{longtable}{p{0.22\linewidth}p{0.34\linewidth}p{0.32\linewidth}}
\toprule
现象 & 可能原因 & 硬件/OS 联系 \\
\midrule
设备利用率高且 await 高 & 队列拥塞或设备慢 & 块层请求和设备服务时间 \\
中断集中在一个 CPU & IRQ affinity 不均 & 多队列和 CPU 拓扑 \\
IOMMU fault & DMA 地址或权限错误 & DMA mapping 生命周期 \\
reset/timeout 日志 & 设备或驱动失去响应 & in-flight 请求恢复 \\
网络 drop 增加 & ring 满或协议栈处理慢 & NAPI、socket backlog、应用读取 \\
\bottomrule
\end{longtable}

这些现象应该能回到前文的状态机：请求是否已经提交，设备是否拥有 buffer，completion 是否回来，等待者是否被唤醒，错误是否传播到用户态。

\subsection{第一阶段负面测试思维}

高质量理解必须包含负面路径。只验证正常读写成功，无法证明边界正确。第一阶段可以用纸上或小实验设计下面这些测试：

\begin{longtable}{p{0.24\linewidth}p{0.34\linewidth}p{0.3\linewidth}}
\toprule
机制 & 应注入的失败 & 期望结果 \\
\midrule
用户指针检查 & 传入无效用户地址 & 返回错误或信号，内核不崩溃 \\
page fault & 访问无 VMA 地址 & 进程失败，不影响其他进程 \\
COW & fork 后父子分别写同一地址 & 内容分离，互不污染 \\
DMA 生命周期 & completion 前释放 buffer 的模拟 & 测试应能发现所有权错误 \\
中断处理 & completion 丢失或迟到 & 超时路径唤醒等待者并报告错误 \\
设备复位 & in-flight 请求未完成时 reset & 请求被 retry 或 fail，不永久睡眠 \\
持久化 & write 后崩溃模拟 & 无 fsync 不承诺持久，有 fsync 按语义恢复 \\
\bottomrule
\end{longtable}

负面测试的价值是逼迫我们说清楚“不变量”。例如“completion 前不能释放 buffer”来自设备仍可能 DMA 的事实；“PTE 改权限后要失效 TLB”关系到隔离条件；“write 不等于持久”来自 page cache、设备 cache 和崩溃恢复共同决定的语义。

\subsection{第一阶段最终验收清单}

下面的清单可以作为第一阶段完成标准。每一项都要求能讲出硬件来源、内核对象、正常路径、失败路径和可观察证据。

\begin{itemize}
\tightlist
\item 能从复位向量讲到内核初始化：固件、内存地图、页表、异常入口、驱动 probe。
\item 能从系统调用讲到特权边界：寄存器 ABI、内核栈、用户指针、返回路径。
\item 能从一次 load/store 讲到 MMU：VMA、PTE、TLB、fault、COW、权限。
\item 能从定时器中断讲到调度：tick、runtime accounting、runqueue、context switch。
\item 能从 DMA 提交讲到 completion：描述符、屏障、IOMMU、中断、等待队列。
\item 能从 file read 讲到硬件设备：fd、page cache、filesystem、block layer、NVMe。
\item 能从网络包讲到 socket：RX ring、NAPI、skb、TCP receive queue、epoll。
\item 能从 write 讲到崩溃恢复：dirty folio、writeback、journal、flush、fsync。
\item 能解释多核带来的成本：cache line、锁竞争、IPI、TLB shootdown、NUMA。
\item 能用至少一种系统证据验证自己的解释，而不是只停留在定义。
\end{itemize}

这份清单完成后，第一阶段才算过关。后续进入进程、内存、文件系统和网络时，读者已经有一条硬件到 OS 的主线，不会把每章当成孤立概念。
